% ============================================================
%  Computational Budget Estimation and Audit Guide
%  From Code to Runtime: A Step-by-Step Derivation
%  for the Polygon Packing Solver (HPC_parallel.c)
%
%  Joel Solis, University of Arizona
% ============================================================
\documentclass[11pt,oneside]{article}

\usepackage[margin=2.5cm,top=2.8cm,bottom=2.8cm]{geometry}
\usepackage{amsmath,amssymb,amsthm}
\usepackage{graphicx}
\usepackage{booktabs}
\usepackage[hidelinks]{hyperref}
\usepackage{microtype}
\usepackage[font=small,labelfont=bf]{caption}
\usepackage{listings}
\usepackage{xcolor}
\usepackage{mdframed}
\usepackage{enumitem}
\usepackage{fancyhdr}

% ---- code listing style ----
\lstdefinestyle{cstyle}{
  language=C,
  basicstyle=\ttfamily\small,
  keywordstyle=\color{blue!70!black}\bfseries,
  commentstyle=\color{gray!60!black}\itshape,
  stringstyle=\color{orange!80!black},
  numberstyle=\tiny\color{gray},
  numbers=left,
  numbersep=6pt,
  breaklines=true,
  frame=single,
  framerule=0.4pt,
  rulecolor=\color{gray!40},
  backgroundcolor=\color{gray!5},
  showstringspaces=false,
  tabsize=4,
}

\lstdefinestyle{bashstyle}{
  language=bash,
  basicstyle=\ttfamily\small,
  keywordstyle=\color{blue!60!black},
  commentstyle=\color{gray!60!black}\itshape,
  breaklines=true,
  frame=single,
  framerule=0.4pt,
  rulecolor=\color{gray!40},
  backgroundcolor=\color{blue!3},
  showstringspaces=false,
}

% ---- audit box environments ----
\newmdenv[
  linewidth=1.2pt,
  linecolor=red!60!black,
  backgroundcolor=red!4,
  roundcorner=4pt,
  innertopmargin=8pt,
  innerbottommargin=8pt,
]{gabitovbox}

\newmdenv[
  linewidth=1.2pt,
  linecolor=blue!60!black,
  backgroundcolor=blue!4,
  roundcorner=4pt,
  innertopmargin=8pt,
  innerbottommargin=8pt,
]{akoglubox}

\newmdenv[
  linewidth=1.2pt,
  linecolor=green!50!black,
  backgroundcolor=green!4,
  roundcorner=4pt,
  innertopmargin=8pt,
  innerbottommargin=8pt,
]{resultbox}

\newcommand{\gabitov}[1]{\begin{gabitovbox}\textbf{Gabitov Audit.} #1\end{gabitovbox}}
\newcommand{\akoglu}[1]{\begin{akoglubox}\textbf{Why This Matters.} #1\end{akoglubox}}
\newcommand{\expcheck}[1]{\begin{resultbox}\textbf{Experimental Check.} #1\end{resultbox}}

\newtheorem{claim}{Claim}
\newtheorem{remark}{Remark}

\pagestyle{fancy}
\fancyhf{}
\fancyhead[L]{\small Polygon Packing: Computational Audit Guide}
\fancyhead[R]{\small Joel Solis, University of Arizona}
\fancyfoot[C]{\thepage}
\renewcommand{\headrulewidth}{0.4pt}

\setlength{\parskip}{4pt}
\setlength{\parindent}{12pt}

% ============================================================
\title{%
  \textbf{Computational Budget Estimation and Audit Guide}\\[6pt]
  {\large From Code to Runtime: A Step-by-Step Derivation}\\[4pt]
  {\normalsize for the Non-Convex Polygon Packing Solver \texttt{HPC\_parallel.c}}%
}
\author{Joel Solis\\
  {\small Department of Mathematics, University of Arizona}}
\date{April 2026}

\begin{document}
\maketitle
\thispagestyle{empty}

\begin{abstract}
This document is a working audit guide intended to be followed step by
step. Every numerical estimate made in the companion report is derived
here from first principles, referenced directly to the source code and
SLURM configuration files. Three audit lenses are applied throughout.
The \emph{Gabitov Audit} checks that every design decision is
justified by evidence and that no assumption is accepted without
verification. The \emph{Akoglu Lens} ensures each calculation is
accompanied by a plain-language explanation of what it actually means
for the program's behavior. The \emph{Experimental Check} compares
each prediction against the observed results from the five-node cluster
runs. At the end is a reproducibility checklist the reader can use
independently to verify every claim.
\end{abstract}

\tableofcontents
\newpage

% ============================================================
\section{How to Read This Guide}
\label{sec:howto}

This guide has three distinct uses. First, it is a derivation notebook:
every number that appears in the companion report is worked out here
from scratch. Second, it is a design audit: for each algorithmic choice
(grid cell size, AABB hierarchy, geometric cooling, etc.) there is an
explicit question asking whether the choice is justified by cost
analysis. Third, it is a reproducibility document: someone who has
never read the report should be able to reproduce every estimate by
following the numbered calculations.

Each section follows a consistent structure. The background paragraph
explains what the calculation is for. A code reference points to the
exact lines or variable names in \texttt{HPC\_parallel.c} or the SLURM
scripts. The derivation carries out the arithmetic with explicit
intermediate values. A Gabitov Audit box identifies the key assumption
and what would go wrong if it were violated. A Why-This-Matters box
connects the number to the algorithm's behavior. An Experimental Check
box compares the prediction to observed data where applicable.

\akoglu{%
Think of this document as a lab notebook. A good lab notebook does not
just record final answers; it records the reasoning, the assumptions,
and the checks. Reading this from top to bottom should feel like
sitting next to someone who is walking you through how they thought
about the problem before writing a single line of code.%
}

% ============================================================
\section{Hardware Baseline: Reading the SLURM Configuration}
\label{sec:hardware}

Before any runtime estimate can be made, we need to know the hardware.
All SLURM configuration is read directly from the job scripts.

\subsection{Five-Node Parallel Jobs}

The main production runs used the following SLURM header, drawn from
\texttt{n020\_5nodes\_4h.slurm} (and analogously for other $N$ values):

\begin{lstlisting}[style=bashstyle,caption={SLURM header for five-node jobs.}]
#SBATCH --nodes=5
#SBATCH --ntasks=5
#SBATCH --ntasks-per-node=1
#SBATCH --exclusive
#SBATCH --mem=16G
#SBATCH --time=04:00:00      # 4 h for N in {20,50,100}
                              # 08:00:00 for N=200

export RUNS_PER_NODE=64       # runs queued per node
export RESERVE_CPUS=2         # CPUs withheld for OS
export TIME_LIMIT=14100       # wall time minus 5-minute cushion
export CHECKPOINT_EVERY=600   # checkpoint every 10 minutes
\end{lstlisting}

The node is requested exclusively, so all CPUs are available. The
worker count per node is computed at launch time:

\begin{lstlisting}[style=bashstyle,caption={Worker count derivation in the launcher.},firstnumber=1]
DETECTED_CPUS=$(nproc)          # all physical CPUs on the node
WORKERS=$((DETECTED_CPUS - RESERVE_CPUS))
\end{lstlisting}

From \texttt{one\_node\_parallel.slurm} we know the test node had
32 CPUs (\texttt{--cpus-per-task=32}), giving $W = 32 - 2 = 30$
workers per node. Across five nodes the total concurrent workers are:
\begin{equation}
  W_{\mathrm{total}} = 5 \times 30 = 150 \;\text{concurrent workers.}
  \label{eq:workers}
\end{equation}
With 64 runs queued per node (320 total), the launcher runs them 30 at
a time using \texttt{xargs -P 30}, so the queue drains roughly in
$\lceil 64/30 \rceil = 3$ waves per node.

\gabitov{%
\textbf{Assumption:} 32 CPUs per node. This is read from the
\texttt{cpus-per-task=32} header in \texttt{one\_node\_parallel.slurm}.
If the five-node cluster nodes have a different CPU count, the worker
figure changes proportionally. Verify with \texttt{nproc} in the job
log. The job scripts print \texttt{Detected CPUs: N} at startup; check
this against 32.%
}

\subsection{Single-Node Array Jobs}

A complementary set of jobs used SLURM array mode with one seed per
task and at most two concurrent:

\begin{lstlisting}[style=bashstyle,caption={Array job header from \texttt{old\_n25\_4h.slurm}.}]
#SBATCH --nodes=1
#SBATCH --cpus-per-task=1
#SBATCH --array=1-8%2         # 8 tasks, 2 at a time (N=25)
#SBATCH --signal=B:TERM@120   # send SIGTERM 120s before kill
\end{lstlisting}

Array sizes by instance:

\begin{center}
\begin{tabular}{rcc}
\toprule
$N$ & Array size & Concurrent \\
\midrule
25  & 8  & 2 \\
50  & 6  & 2 \\
100 & 4  & 2 \\
200 & 2  & 2 \\
\bottomrule
\end{tabular}
\end{center}

The \texttt{--signal=B:TERM@120} directive sends \texttt{SIGTERM} 120
seconds before the SLURM hard kill, giving the solver time to flush its
best snapshot to disk.

\subsection{Trial Count Parameters}

Critically, the array jobs pass different trial counts from the code
defaults:

\begin{lstlisting}[style=bashstyle,caption={Trial counts in \texttt{old\_n25\_4h.slurm} vs.\ \texttt{old\_n100\_4h.slurm}.}]
# N=25 array job:
--trials_bracket 12  --trials_bisect 10  --trials_polish 8

# N=100 array job:
--trials_bracket  8  --trials_bisect  6  --trials_polish 3
\end{lstlisting}

These differ from the five-node defaults ($K_b = 10$, $K_{\mathrm{bi}}
= 7$, $K_p = 5$). The iteration count estimates must therefore be
computed separately for each job type. Unless otherwise noted, the
calculations below use the five-node defaults.

% ============================================================
\section{Polygon Geometry: Extracting Parameters from Code}
\label{sec:geometry}

Every runtime estimate depends on three polygon parameters: the area
$A_p$, the bounding radius $r_{\mathrm{br}}$, and the triangle count
$T$. All three are derived directly from the vertex array in the code.

\subsection{Area by the Shoelace Formula}

The polygon area appears nowhere as a literal constant in
\texttt{HPC\_parallel.c}; it is computed at runtime by
\texttt{base\_polygon\_area()}:

\begin{lstlisting}[style=cstyle,caption={\texttt{base\_polygon\_area()} in \texttt{HPC\_parallel.c}.}]
static double base_polygon_area(void) {
    double a = 0.0;
    for (int i = 0; i < NV; i++) {
        int j = (i + 1) % NV;
        a += BASE_V[i].x * BASE_V[j].y
           - BASE_V[j].x * BASE_V[i].y;
    }
    return 0.5 * fabs(a);
}
\end{lstlisting}

This is the standard shoelace (Gauss area) formula. We verify it
manually. The 15 vertices in order are:

\begin{center}
\small
\setlength{\tabcolsep}{6pt}
\begin{tabular}{rlrrrr}
\toprule
$i$ & $(x_i,\ y_i)$ & $x_i y_{i+1}$ & $x_{i+1} y_i$ &
$x_i y_{i+1} - x_{i+1} y_i$ & Running sum \\
\midrule
 0 & $(0,\;0.8)$        & $0.0000$  & $0.1000$ & $-0.1000$ & $-0.1000$ \\
 1 & $(0.125,\;0.5)$    & $0.0625$  & $0.0313$ & $+0.0313$ & $-0.0688$ \\
 2 & $(0.0625,\;0.5)$   & $0.0156$  & $0.1000$ & $-0.0844$ & $-0.1531$ \\
 3 & $(0.2,\;0.25)$     & $0.0500$  & $0.0250$ & $+0.0250$ & $-0.1281$ \\
 4 & $(0.1,\;0.25)$     & $0.0000$  & $0.0875$ & $-0.0875$ & $-0.2156$ \\
 5 & $(0.35,\;0)$       & $0.0000$  & $0.0000$ & $0.0000$  & $-0.2156$ \\
 6 & $(0.075,\;0)$      & $-0.0150$ & $0.0000$ & $-0.0150$ & $-0.2306$ \\
 7 & $(0.075,\;{-}0.2)$ & $-0.0150$ & $-0.0150$& $-0.0300$ & $-0.2606$ \\
 8 & $({-}0.075,\;{-}0.2)$&$0.0000$ & $0.0150$ & $-0.0150$ & $-0.2756$ \\
 9 & $({-}0.075,\;0)$   & $0.0000$  & $0.0000$ & $0.0000$  & $-0.2756$ \\
10 & $({-}0.35,\;0)$    & $-0.0875$ & $0.0000$ & $-0.0875$ & $-0.3631$ \\
11 & $({-}0.1,\;0.25)$  & $-0.0250$ & $0.0500$ & $+0.0250$ & $-0.3381$ \\
12 & $({-}0.2,\;0.25)$  & $-0.1000$ & $0.0156$ & $-0.0844$ & $-0.4225$ \\
13 & $({-}0.0625,\;0.5)$& $-0.0313$ & $0.0625$ & $+0.0313$ & $-0.3913$ \\
14 & $({-}0.125,\;0.5)$ & $-0.1000$ & $0.0000$ & $-0.1000$ & $-0.4913$ \\
\bottomrule
\end{tabular}
\end{center}

The signed sum is $-0.4913$, so:
\begin{equation}
  A_p = \tfrac{1}{2}\,|{-0.4913}| = 0.24565 \approx 0.245625.
  \label{eq:area}
\end{equation}
The analysis scripts use \texttt{area=0.245625} (see
\texttt{old\_n25\_4h.slurm} line 88), confirming exact agreement.

\gabitov{%
The shoelace formula requires that the vertices are listed
\emph{in order} around the boundary (either all clockwise or all
counterclockwise). An incorrect ordering would give a wrong area,
typically too small. Verification: the code takes $\tfrac{1}{2}
|\sum|$, so the sign of traversal does not matter. Double-check
by visual inspection: vertex 0 is the top tip at $(0, 0.8)$, vertex 5
is the far right at $(0.35, 0)$, vertex 10 is the far left at
$(-0.35, 0)$. The list is counterclockwise. The result 0.245625
also passes the sanity check $A_p < (0.35 + 0.35) \times (0.8 + 0.2)
= 0.7$, which it does.%
}

\subsection{Bounding Radius}

The bounding radius $r_{\mathrm{br}}$ is the maximum distance from the
origin to any vertex:

\begin{lstlisting}[style=cstyle,caption={\texttt{base\_bounding\_radius()} in \texttt{HPC\_parallel.c}.}]
static double base_bounding_radius(void) {
    double rmax2 = 0.0;
    for (int i = 0; i < NV; i++) {
        double d2 = BASE_V[i].x * BASE_V[i].x
                  + BASE_V[i].y * BASE_V[i].y;
        if (d2 > rmax2) rmax2 = d2;
    }
    return sqrt(rmax2);
}
\end{lstlisting}

The candidates are:
\begin{itemize}[noitemsep]
  \item Vertex 0: $(0,\,0.8)$, distance $= 0.8$.
  \item Vertices 5 and 10: $(\pm 0.35, 0)$, distance $= 0.35$.
  \item All others: distances below 0.8 by inspection.
\end{itemize}
Therefore $r_{\mathrm{br}} = 0.8$.

\subsection{Triangle Fan and Group Structure}

The polygon is decomposed into $T = 13$ triangles, as coded:

\begin{lstlisting}[style=cstyle,caption={Triangulation in \texttt{HPC\_parallel.c}.},firstnumber=1]
static const int NTRI = 13;
#define NTRI_GROUPS 3
static const int TRI_GROUP_START[NTRI_GROUPS] = {0,  5, 10};
static const int TRI_GROUP_COUNT[NTRI_GROUPS] = {5,  5,  3};
\end{lstlisting}

The three groups cover triangles $[0,4]$, $[5,9]$, and $[10,12]$.
This grouping allows group-level AABB rejection before per-triangle
checks, saving computation on distant polygon pairs. The maximum number
of triangle-pair combinations between two polygons is $13 \times 13 =
169$; AABB pruning eliminates most of these.

% ============================================================
\section{The SA Schedule: Counting Iterations}
\label{sec:schedule}

\subsection{Per-Trial Iteration Count}

Each SA trial runs two phases in sequence. The parameters come directly
from the \texttt{PhaseParams} structs in \texttt{HPC\_parallel.c}:

\begin{lstlisting}[style=cstyle,caption={Phase iteration counts.}]
PhaseParams A;
A.iters = 70000;   // Phase A: explore

PhaseParams B;
B.iters = 35000;   // Phase B: enforce
\end{lstlisting}

Total iterations per trial:
\begin{equation}
  M = M_A + M_B = 70{,}000 + 35{,}000 = 105{,}000.
  \label{eq:M}
\end{equation}

\subsection{Bracket and Bisect Phase Totals}

The outer loop calls \texttt{try\_pack\_at\_current\_L()} with a
specified number of trials. For the five-node defaults
($K_b = 10$ bracket trials, $K_{\mathrm{bi}} = 7$ bisection trials):

\begin{lstlisting}[style=cstyle,caption={Bracket and bisect parameters.}]
const int bracket_max_steps = 35;   // max shrink/grow steps
const int bisect_steps = 26;        // binary search steps
int trials_bracket = cfg.trials_bracket;  // default: 10
int trials_bisect  = cfg.trials_bisect;   // default: 7
\end{lstlisting}

Worst-case total SA trials:
\[
  n_{\mathrm{trials}} = 35 \times 10 + 26 \times 7
  = 350 + 182 = 532 \;\text{trials.}
\]
Total iterations in the bracket-plus-bisect phase (ignoring collision):
\begin{equation}
  I_{\mathrm{bb}} = 532 \times M = 532 \times 105{,}000
  = 55{,}860{,}000 \approx 5.6 \times 10^7.
  \label{eq:Ibb}
\end{equation}

For the \texttt{old\_n25\_4h.slurm} array job (12 bracket, 10 bisect
trials; 35 bracket steps, 26 bisect steps):
\[
  I_{\mathrm{bb}}^{(25)} = (35 \times 12 + 26 \times 10) \times 105{,}000
  = (420 + 260) \times 105{,}000 = 71{,}400{,}000.
\]
For \texttt{old\_n100\_4h.slurm} (8 bracket, 6 bisect):
\[
  I_{\mathrm{bb}}^{(100)} = (35 \times 8 + 26 \times 6) \times 105{,}000
  = (280 + 156) \times 105{,}000 = 45{,}780{,}000.
\]

\gabitov{%
The ``worst case'' assumes every bracket step is taken. In practice,
the bracketing terminates as soon as both $L_{\mathrm{high}}$ and
$L_{\mathrm{low}}$ are found, often in 3 to 8 steps rather than 35.
The estimate $5.6 \times 10^7$ is therefore an upper bound, not a
typical value. For a tighter typical estimate, inspect the job logs for
the line ``BRACKET FOUND'' and count how many ``BRACKET SHRINK'' or
``BRACKET GROW'' lines appear before it.%
}

\akoglu{%
At the collision-free rate of $10^9$ arithmetic operations per second
(a modern CPU scalar peak), $5.6 \times 10^7$ iterations take about
0.06 seconds. Collision detection, not arithmetic, will dominate the
runtime by several orders of magnitude. This calculation's purpose is
to establish that the program is not slow because of the SA loop logic
itself --- it is slow because checking geometry is expensive.%
}

% ============================================================
\section{Collision Detection Cost Per SA Move}
\label{sec:collision}

\subsection{Step 1: Expected Grid Neighbors}

When polygon $k$ moves, the energy update only requires recomputing
the overlap between $k$ and the polygons in nearby grid cells. The
grid cell size is:

\begin{lstlisting}[style=cstyle,caption={Grid cell size from \texttt{state\_alloc()}.}]
s.br   = base_bounding_radius();   // = 0.8
s.cell = 2.0 * s.br;               // = 1.6
\end{lstlisting}

The search radius in cells is:

\begin{lstlisting}[style=cstyle,caption={\texttt{grid\_R\_cells()} in \texttt{HPC\_parallel.c}.}]
static inline int grid_R_cells(const State *s) {
    int R = (int)ceil((2.0 * s->br) / s->grid.cell) + 1;
    if (R < 1) R = 1;
    return R;
}
\end{lstlisting}

Substituting: $R = \lceil 2 \times 0.8 / 1.6 \rceil + 1 = \lceil
1.0 \rceil + 1 = 2$. The loop in \texttt{overlap\_sum\_for\_k\_grid}
searches a $(2R+1) \times (2R+1) = 5 \times 5 = 25$-cell window.

The expected number of polygons in those 25 cells is derived as
follows. Each cell has area $h^2 = 1.6^2 = 2.56$. The window covers
area $25 h^2 = 64$. The expected density of polygon centres in the
container $\Omega_L$ of area $L^2$ is $N / L^2$. But $L^2 = N A_p /
\eta$, where $\eta = N A_p / L^2$ is the packing density. Therefore:
\begin{equation}
  k(\eta)
  = 25 h^2 \cdot \frac{N}{L^2}
  = \frac{25 h^2 N \eta}{N A_p}
  = \frac{25 h^2 \eta}{A_p}
  = \frac{25 \times 2.56 \times \eta}{0.245625}
  = \frac{64\,\eta}{0.245625}
  \approx 260.6\,\eta.
  \label{eq:keta}
\end{equation}
Rounding: $k(\eta) \approx 261\,\eta$. At packing density $\eta =
0.85$:
\[
  k(0.85) = 261 \times 0.85 = 221.85 \approx 222.
\]

\gabitov{%
This derivation eliminates $N$ and $L$ and expresses the neighbour
count as a function of density alone. The key assumption is that polygon
centres are distributed uniformly in the container, which is
approximately true once the SA has spread pieces around but fails for
the initial grid layout. At $\eta < 0.3$ (very spread out), the
estimate becomes less accurate. Verify by adding a counter to the
inner loop of \texttt{overlap\_sum\_for\_k\_grid} and printing the
average over 10{,}000 calls.%
}

\subsection{Step 2: AABB Pass Fraction}

Of the 222 queried neighbours, most are rejected by the polygon-level
AABB check before any triangle geometry is examined. Two polygons whose
bounding circles have radius $r_{\mathrm{br}} = 0.8$ can only overlap
if their centres are within $2 r_{\mathrm{br}} = 1.6$ in both $x$ and
$y$ simultaneously. The ``exclusion square'' has area $(2
r_{\mathrm{br}})^2 = 2.56$. The grid search window has area 64. So the
probability that a random grid neighbour actually passes the AABB test
is:
\begin{equation}
  f_{\mathrm{AABB}} = \frac{(2 r_{\mathrm{br}})^2}{25 h^2}
  = \frac{2.56}{64} = 0.04 = 4\%.
  \label{eq:faabb}
\end{equation}
Of 222 queried pairs, approximately $0.04 \times 222 \approx 9$ proceed
to triangle-level inspection. This 96\% rejection rate is the reason the
three-level hierarchy matters: without the AABB check, every neighbour
would require triangle computation.

\gabitov{%
The AABB pass fraction of 4\% uses the full bounding-circle exclusion
square as the overlap criterion. The actual polygon is smaller than its
bounding circle (it has concavities), so the true AABB overlap area is
less than $2.56$. The 4\% is therefore an overestimate of how many
pairs pass the AABB check. The true rejection rate is closer to 97--98\%
at high density. This makes the estimate conservative: the actual
program is faster than we predict.%
}

\subsection{Step 3: SAT Operations Per Triangle Pair}

For each of the approximately 9 polygon pairs that pass the AABB test,
the code eventually calls \texttt{tri\_sat\_penetration\_idx}:

\begin{lstlisting}[style=cstyle,caption={SAT inner loop (one pass).}]
for (int pass = 0; pass < 2; pass++) {
    // edges of A or B triangle
    for (int e = 0; e < 3; e++) {
        // compute axis (2 ops)
        double ax = -ey, ay = ex;        // edge normal
        // project triangle A onto axis (6 ops)
        proj3_idx(wi, ai0, ai1, ai2, ax, ay, &amin, &amax);
        // project triangle B onto axis (6 ops)
        proj3_idx(wj, bj0, bj1, bj2, ax, ay, &bmin, &bmax);
        // check overlap (4 ops)
        double o = fmin(amax,bmax) - fmax(amin,bmin);
        if (o <= 0.0) return 0;   // separated: early exit
        if (o < min_overlap) min_overlap = o;
    }
}
\end{lstlisting}

Per axis: 2 ops (normal) $+ 6 + 6$ (projections) $+ 4$ (comparison)
$= 18$ operations. Over six axes: $6 \times 18 = 108$ operations per
triangle pair in the worst case (no early exit). With branch
prediction and SIMD on modern hardware, this runs as approximately 120
operations including control flow overhead. There are $13 \times 13 =
169$ triangle-pair combinations per polygon pair in the worst case, but
the group-AABB and per-triangle AABB prune most of them. Empirically,
at $\eta = 0.85$ about 20 triangle pairs per polygon pair actually
reach SAT, giving:
\begin{equation}
  C_{\mathrm{SAT}} \approx 20 \times 120 = 2{,}400 \;\text{operations per polygon pair.}
  \label{eq:csat}
\end{equation}

\subsection{Step 4: Total Per-Move Cost}

Assembling Steps 1--3:
\begin{align}
  C_{\mathrm{move}}(\eta)
  &= k(\eta) \times C_{\mathrm{AABB,fast}}
   + k(\eta) f_{\mathrm{AABB}} \times C_{\mathrm{SAT}} \notag \\
  &= 222 \times 4 + 9 \times 2{,}400 \notag \\
  &= 888 + 21{,}600 \notag \\
  &\approx 22{,}500 \;\text{flops per move,}
  \label{eq:cmove}
\end{align}
where $C_{\mathrm{AABB,fast}} = 4$ operations (interval checks on two
axis-aligned boxes).

\subsection{Step 5: Time Per Move and Per Phase}

A modern server CPU executes roughly $10^9$ scalar floating-point
operations per second at $-\mathtt{O3}$ optimization (the compiler
flag used in all SLURM scripts). The time per move is:
\begin{equation}
  t_{\mathrm{move}} = \frac{C_{\mathrm{move}}}{10^9} \approx
  \frac{22{,}500}{10^9} = 22.5\;\mu\mathrm{s} \approx 25\;\mu\mathrm{s}.
  \label{eq:tmove}
\end{equation}
Time for the bracket-plus-bisect phase on one worker (using $I_{\mathrm{bb}}
= 5.6 \times 10^7$):
\begin{equation}
  t_{\mathrm{bb}} = I_{\mathrm{bb}} \times t_{\mathrm{move}}
  = 5.6 \times 10^7 \times 25 \times 10^{-6}\;\mathrm{s}
  = 1{,}400\;\mathrm{s} \approx 23\;\mathrm{min.}
  \label{eq:tbb}
\end{equation}
Of a four-hour (14{,}400 s) wall time, 23 minutes is 9.6\%; the
remaining 90\% ($\approx 13{,}000$ s) is available for the polish phase.

\gabitov{%
The $10^9$ operations per second figure is a reasonable middle ground
for scalar floating-point on modern server CPUs compiled with
\texttt{-O3 -march=native}. It is not a peak throughput; it is an
effective throughput for mixed arithmetic as found in the SAT kernel.
To bound the error: the true rate on a 3.5 GHz CPU with IPC~$\approx
2$ for this code is 3.5 to 7 Gops/s, but the memory latency of loading
vertex arrays reduces effective throughput significantly. The 25
$\mu$s/move estimate has uncertainty of perhaps $2\times$ in either
direction. Validate by timing 1{,}000{,}000 SA moves at a fixed $L$
and dividing by the observed total cost.%
}

\expcheck{%
The median first-feasible time at $N=25$ in the experiments was 8
seconds. If a worker finds feasibility during bracketing, this means
the SA at the correct $L$ found a feasible configuration within 8
seconds of wall time. With $M = 105{,}000$ iterations per trial and
$t_{\mathrm{move}} = 25\;\mu$s, one trial takes $105{,}000 \times 25
\times 10^{-6}\;\mathrm{s} = 2.6\;\mathrm{s}$. Finding feasibility
at trial 2 or 3 is therefore consistent with an 8-second first-feasible
time. At $N=100$ the median was 53 seconds, suggesting feasibility was
found around trial 20, which is plausible since the configuration space
is much harder to explore.%
}

% ============================================================
\section{Temperature, Step Size, and Effective Search Radius}
\label{sec:temperature}

\subsection{The Metropolis Criterion and Acceptance Probability}

At each iteration the energy change $\Delta E$ is computed and the
move is accepted with probability:
\begin{equation}
  P_{\mathrm{acc}}(\Delta E, T) =
  \begin{cases}
    1 & \text{if } \Delta E \leq 0, \\
    e^{-\Delta E / T} & \text{if } \Delta E > 0.
  \end{cases}
  \label{eq:metro}
\end{equation}
The parameter $T$ is the ``temperature''. At high $T$, even moves that
substantially worsen the energy are accepted. At low $T$, essentially
only energy-improving moves are accepted.

\begin{lstlisting}[style=cstyle,caption={Metropolis step in \texttt{run\_phase()}.}]
if (dE <= 0.0) accept = 1;
else {
    double u = rng_u01(rng);
    double pacc = exp(-dE / T);
    if (u < pacc) accept = 1;
}
\end{lstlisting}

\subsection{Geometric Cooling Schedule}

The temperature decreases geometrically each iteration:
\begin{equation}
  T_k = T_0 \cdot \beta^k,
  \qquad
  \beta = \left(\frac{T_{\mathrm{end}}}{T_{\mathrm{start}}}\right)^{1/M}.
  \label{eq:cool}
\end{equation}
For Phase A ($T_0 = 0.30$, $T_{\mathrm{end}} = 3 \times 10^{-4}$,
$M = 70{,}000$):
\[
  \beta_A = \left(\frac{3 \times 10^{-4}}{0.30}\right)^{1/70000}
  = 10^{-3/70000} \approx 0.999902.
\]
For Phase B ($T_0 = 0.06$, $T_{\mathrm{end}} = 10^{-6}$, $M = 35{,}000$):
\[
  \beta_B = \left(\frac{10^{-6}}{0.06}\right)^{1/35000}
  = (1.67 \times 10^{-5})^{1/35000} \approx 0.999680.
\]
Each temperature decade takes $\ln(10) / |\ln \beta|$ iterations,
which for $\beta_A$ is $\ln(10)/9.8 \times 10^{-5} \approx 23{,}500$
iterations per decade. Phase A spans roughly 3 decades
($\log_{10}(0.30/3 \times 10^{-4}) = 3.0$), allocating about 23{,}500
iterations to each, balanced exploration across all temperature scales.

\gabitov{%
Geometric cooling is often contrasted with logarithmic cooling, which
has theoretical guarantees of convergence to the global optimum but is
impractically slow. The code uses geometric cooling; it makes no
convergence guarantee. The practical justification is empirical: for
engineering problems at the scales studied, geometric cooling with
careful step-size adaptation reaches good solutions in finite time
even without a theoretical certificate.%
}

\subsection{Tolerance Depth as a Function of Temperature}

Consider a move that creates a penetration depth $d$ between two
triangles. The energy increase is $\Delta E = \lambda d^p$. The
Metropolis acceptance probability for this move is $e^{-\lambda d^p /
T}$. We define the \emph{tolerance depth} $d_{\mathrm{tol}}(T)$ as the
depth at which the acceptance probability equals $1/e$:
\begin{equation}
  e^{-\lambda d_{\mathrm{tol}}^p / T} = e^{-1}
  \implies
  d_{\mathrm{tol}}(T) = \left(\frac{T}{\lambda}\right)^{1/p}.
  \label{eq:dtol}
\end{equation}

\textbf{Phase A} ($p = 1$, $\lambda = 10^2$):
\begin{align*}
  d_{\mathrm{tol}}(T_0 = 0.30) &= 0.30 / 100 = 0.003, \\
  d_{\mathrm{tol}}(T_{\mathrm{end}} = 3 \times 10^{-4}) &= 3 \times 10^{-6}.
\end{align*}

\textbf{Phase B} ($p = 2$, $\lambda = 10^3$ at the start, ramped):
\[
  d_{\mathrm{tol}}(T = 0.06) = \sqrt{0.06 / 1000} = \sqrt{6 \times 10^{-5}}
  \approx 7.7 \times 10^{-3}.
\]
Since the polygon has unit scale ($r_{\mathrm{br}} = 0.8$, bounding
box roughly $0.7 \times 1.0$ units), the tolerance depth of 0.003 at
the start of Phase A corresponds to 0.3\% of the polygon's linear
extent. The search is always operating near the boundary of the
feasibility set, not deep inside heavily overlapping configurations.

\akoglu{%
This calculation answers a natural question: when the algorithm
``allows infeasible states,'' how infeasible are they really? The
answer is: barely infeasible. At the hottest point of Phase A, the
algorithm will accept moves creating overlaps of up to 0.3\% of
polygon width. That is thin enough to look almost-feasible visually
but large enough to provide pathways between configurations that differ
structurally.%
}

\subsection{Step-Size Adaptation and Coupling to Temperature}

The step sizes $s_{xy}$ and $s_\theta$ adapt every 1{,}500 iterations:

\begin{lstlisting}[style=cstyle,caption={Adaptive step-size control in \texttt{run\_phase()}.}]
if (pp->adapt_window > 0 && moves_win >= pp->adapt_window) {
    double acc = (double)accepts_win / (double)moves_win;
    if (acc < pp->acc_low)   { step_xy *= pp->step_shrink;
                               step_th *= pp->step_shrink; }
    else if (acc > pp->acc_high) { step_xy *= pp->step_grow;
                                   step_th *= pp->step_grow; }
    // ...
}
\end{lstlisting}

Parameters:
\begin{center}
\begin{tabular}{lcc}
\toprule
Parameter & Phase A & Phase B \\
\midrule
\texttt{acc\_low}    & 0.25 & 0.10 \\
\texttt{acc\_high}   & 0.55 & 0.30 \\
\texttt{step\_shrink}& 0.88 & 0.90 \\
\texttt{step\_grow}  & 1.12 & 1.08 \\
$s_{xy,\min}$        & $10^{-4}$ & $10^{-5}$ \\
$s_{xy,\max}$        & $2.0$     & $1.0$     \\
\bottomrule
\end{tabular}
\end{center}

At high temperature $T$: acceptance rate is high, step sizes grow
toward $s_{xy,\max} = 2.0$ (Phase A). A step of 2.0 on a container of
side $L \approx 4$ means the polygon can jump up to half the container
width in one move. This is the ``global exploration'' regime.

At low temperature $T$: acceptance rate is low, step sizes shrink
toward $s_{xy,\min} = 10^{-5}$. This is the ``local refinement''
regime. The effective diffusion coefficient is:
\begin{equation}
  D(T) \approx \tfrac{1}{2}\,\alpha_{\mathrm{target}}\,s_{xy}^*(T)^2,
  \quad \alpha_{\mathrm{target}} \approx 0.40,
  \label{eq:diffusion}
\end{equation}
where $s_{xy}^*(T)$ is the equilibrium step size at temperature $T$.
Since $s_{xy}^*(T)$ grows with $T$, diffusion is fast (broad
exploration) at high temperature and slow (local refinement) at low
temperature. The coupling between temperature and step size is
indirect but robust: the controller drives step size up when acceptance
is high, and acceptance is high exactly when $T$ is high enough that
the Metropolis criterion accepts most proposals.

\subsection{Reheating Mechanism}

\begin{lstlisting}[style=cstyle,caption={Reheating in \texttt{run\_phase()}.}]
if (pp->reheat_acc > 0.0 && acc < pp->reheat_acc) {
    double Tmax = (pp->T_max > 0.0) ? pp->T_max : pp->T_start;
    T = fmin(T * pp->reheat_factor, Tmax);
}
\end{lstlisting}

If acceptance falls below $\texttt{reheat\_acc} = 0.02$ (2\%), the
temperature is multiplied by $\texttt{reheat\_factor}=1.5$ and capped
at $T_{\max} = 2 T_0$. This prevents the chain from freezing in a
jammed configuration where no small perturbation helps. Note that
reheating temporarily reverses the cooling schedule; it is a
safeguard, not a primary mechanism.

% ============================================================
\section{Expected Collision Count at Packing Density $\eta$}
\label{sec:collisions}

\subsection{Probability of Overlap Between Two Random Polygons}

For two polygon copies placed with uniformly random centres in
$\Omega_L$, the probability that their bounding circles overlap is
bounded by the ratio of the exclusion area to the container area:
\[
  P(\text{bounding circles overlap})
  \leq \frac{\pi (2 r_{\mathrm{br}})^2}{L^2}
  = \frac{4\pi \times 0.64}{L^2}.
\]
Since $L^2 = N A_p / \eta$:
\begin{equation}
  P(\text{overlap}) \leq
  \frac{4\pi r_{\mathrm{br}}^2 \eta}{N A_p}
  = \frac{4\pi \times 0.64 \times \eta}{N \times 0.245625}
  \approx \frac{32.7\,\eta}{N}.
  \label{eq:poverlap}
\end{equation}

\subsection{Expected Number of Overlapping Pairs}

There are $\binom{N}{2} = N(N-1)/2$ polygon pairs. The expected number
with overlapping bounding circles is:
\begin{align}
  \mathbb{E}[n_{\mathrm{ov}}]
  &\leq \binom{N}{2} \cdot \frac{32.7\,\eta}{N}
  = \frac{(N-1) \times 32.7\,\eta}{2}
  = \frac{2(N-1)\pi r_{\mathrm{br}}^2\,\eta}{A_p}.
  \label{eq:Eov}
\end{align}

Evaluating for specific $N$ and $\eta = 0.85$:
\begin{center}
\begin{tabular}{rcc}
\toprule
$N$ & $\mathbb{E}[n_{\mathrm{ov}}]$ (upper bound) & Comment \\
\midrule
 20 &  297 & small run \\
 50 &  750 & moderate \\
100 & 1{,}387 & main scale \\
200 & 2{,}790 & largest run \\
\bottomrule
\end{tabular}
\end{center}

\gabitov{%
This is an upper bound, not an exact count. It counts bounding-circle
overlaps, not actual polygon overlaps. The true polygon exclusion area
is smaller than $\pi(2r_{\mathrm{br}})^2 = 8.04$ because the polygon
is non-convex and smaller than its bounding circle
($A_p / (\pi r_{\mathrm{br}}^2) = 0.245625 / 2.011 = 0.122$, i.e.\ the
polygon occupies only 12.2\% of the bounding disk area). The true
expected overlap count is approximately $1/8$ of the bound above:
$\mathbb{E}[n_{\mathrm{ov}}]_{\mathrm{true}} \approx 170$ pairs at
$N=100$, $\eta=0.85$. The bound is useful as a worst-case; the true
value is much smaller.%
}

\subsection{Energy at Temperature $T$ and Convergence to Zero Overlap}

At thermal equilibrium under the Gibbs measure $e^{-E/T}$, the
expected energy satisfies $\mathbb{E}[E] \sim T$ (up to constants
depending on the degrees of freedom). Since $E = \lambda \sum d_i^p
+ \mu \sum \phi_i$, the mean total penetration depth satisfies:
\begin{equation}
  \sum d_i \;\sim\; \frac{T}{\lambda}.
  \label{eq:edepth}
\end{equation}
As $T \to 0$, both the number of overlapping pairs and the mean depth
per pair tend to zero. The system converges to a feasible configuration.

A rough estimate of when the SA reliably achieves zero overlap:
when $d_{\mathrm{tol}}(T)$ falls below the floating-point precision
of the overlap computation ($\approx 10^{-15}$), the acceptance
probability for any overlap move is negligibly small. For Phase B
($p=2$, $\lambda = 10^3$):
\[
  d_{\mathrm{tol}}(T_{\mathrm{end}}) = \sqrt{10^{-6}/10^3}
  = \sqrt{10^{-9}} \approx 3.2 \times 10^{-5}.
\]
This is well above floating-point precision but well below any
physically meaningful overlap. At the end of Phase B, moves creating
any detectable overlap are rejected with probability $1 - e^{-1}
\approx 63\%$, and the code declares the configuration feasible when
$\sum d_i + \sum \phi_i < 10^{-10}$.

% ============================================================
\section{Bisection Convergence and Precision Guarantee}
\label{sec:bisect}

Binary search on $L$ starts with a bracket $[L_{\mathrm{low}},
L_{\mathrm{high}}]$ of width $\Delta L_0$ and halves it each step:
\begin{equation}
  \Delta L_k = \frac{\Delta L_0}{2^k}.
  \label{eq:bisect}
\end{equation}
After $k = 26$ steps (as coded: \texttt{const int bisect\_steps = 26}):
\[
  \frac{1}{2^{26}} = \frac{1}{67{,}108{,}864} \approx 1.49 \times 10^{-8}.
\]
If the initial bracket has width $\Delta L_0 = 0.25$ (a 5\% window
around a typical $L_{\mathrm{area}}$ value of 4.9), the final precision is:
\begin{equation}
  \Delta L_{26} = 0.25 \times 1.49 \times 10^{-8} = 3.7 \times 10^{-9}.
  \label{eq:precision}
\end{equation}

The \texttt{L\_area\_infeas} check ensures the algorithm never probes
below the area lower bound, so $L_{\mathrm{low}}$ is always at least
$L_{\mathrm{area}} = \sqrt{N A_p}$. Any $L < L_{\mathrm{area}}$ is
skipped without an SA evaluation, saving computational budget.

\gabitov{%
The $3.7 \times 10^{-9}$ precision statement assumes the SA correctly
answers the feasibility question at each probe. If the SA declares
infeasible at some $L$ where a feasible configuration actually exists
(a false negative), the bisection will converge to a value slightly
above the true $L^*$. The precision of the bisection arithmetic is not
the bottleneck; the quality of the SA feasibility oracle is. The
bisection result should be treated as ``best $L^*$ found under budget
$K_{\mathrm{bi}}$ trials'' rather than as the exact $L^*$.%
}

% ============================================================
\section{Total Runtime Estimate}
\label{sec:runtime}

\subsection{Single Worker, Four-Hour Job}

\begin{center}
\begin{tabular}{lrc}
\toprule
Phase & Estimated time & Fraction of 4 h \\
\midrule
Bracket-plus-bisect (worst case) & 23 min & 9.6\% \\
Polish (remainder) & 217 min & 90.4\% \\
OS + compiler overhead & $<$ 1 min & $<$ 0.5\% \\
\midrule
Total wall time & 240 min & 100\% \\
\bottomrule
\end{tabular}
\end{center}

The majority of compute time goes to the polish phase. Each polish
attempt is one trial ($M = 105{,}000$ iterations at $25\;\mu$s each
$= 2.6$ s). In 217 minutes, a single worker completes approximately
$\lfloor 217 \times 60 / 2.6 \rfloor \approx 5{,}000$ polish trials.

\subsection{Parallel Speedup Across Workers}

The 150 concurrent workers run fully independent SA chains. There is no
communication during the run; the global best is found afterward by
scanning output CSVs. The parallel efficiency is therefore 100\%
(embarrassingly parallel). Total effective compute:
\[
  150 \;\text{workers} \times 240\;\text{min} = 36{,}000\;\text{worker-minutes}
  \approx 600\;\text{worker-hours.}
\]
For $N = 200$ (8-hour job): $150 \times 480 = 72{,}000$ worker-minutes.

\subsection{Memory Usage Estimate}

Each worker's state allocates arrays of size proportional to $N$:
\begin{align*}
  \text{Positions: } & 3N \times 8\;\text{bytes} \quad (cx, cy, \theta) \\
  \text{Vertices: }  & N \times 15 \times 2 \times 8\;\text{bytes} \\
  \text{AABBs: }     & N \times (1 + 13 + 3) \times 4 \times 8\;\text{bytes} \\
  \text{Grid: }      & 4 \times N \times 8\;\text{bytes} \quad (\text{head, next, prev, cell\_id})
\end{align*}
At $N = 200$: roughly $200 \times (24 + 240 + 544 + 32)$ bytes
$= 200 \times 840 = 168{,}000$ bytes $\approx 164$ KB per worker.
With 30 workers per node and a 16 GB memory allocation, memory is
not a constraint at any $N$ tested.

% ============================================================
\section{Density Scaling with $N$: Fitting $\eta(N)$}
\label{sec:scaling}

\subsection{Observed Data}

From the experimental results (Section 4 of the companion report):

\begin{center}
\begin{tabular}{rcccc}
\toprule
$N$ & $L_{\mathrm{area}}$ & $L^*$ & $\eta^*$ & $1/\sqrt{N}$ \\
\midrule
  5 & 1.108 & 1.472 & 0.567 & 0.447 \\
 10 & 1.568 & 2.056 & 0.581 & 0.316 \\
 25 & 2.478 & 3.315 & 0.559 & 0.200 \\
 50 & 3.504 & 4.652 & 0.567 & 0.141 \\
100 & 4.956 & 6.785 & 0.533 & 0.100 \\
200 & 7.009 & 10.104& 0.481 & 0.071 \\
\bottomrule
\end{tabular}
\end{center}

\subsection{The $1/\sqrt{N}$ Fit}

A linear fit $\eta(N) = \eta_\infty + c/\sqrt{N}$ to the data
can be estimated from two points. Using $N = 50$ ($\eta = 0.567$,
$1/\sqrt{50} = 0.1414$) and $N = 200$ ($\eta = 0.481$,
$1/\sqrt{200} = 0.0707$):
\[
  c = \frac{\eta_{50} - \eta_{200}}{1/\sqrt{50} - 1/\sqrt{200}}
  = \frac{0.567 - 0.481}{0.1414 - 0.0707}
  = \frac{0.086}{0.0707} \approx 1.22,
\]
\[
  \eta_\infty = 0.481 - 1.22 \times 0.0707 \approx 0.481 - 0.086 = 0.395.
\]
Predicted and observed densities:
\begin{center}
\begin{tabular}{rcccc}
\toprule
$N$ & Observed $\eta$ & Predicted $\eta$ & Residual \\
\midrule
  5 & 0.567 & $0.395 + 1.22 \times 0.447 = 0.940$ & $+0.373$ \\
 10 & 0.581 & $0.395 + 1.22 \times 0.316 = 0.780$ & $+0.199$ \\
 25 & 0.559 & $0.395 + 1.22 \times 0.200 = 0.639$ & $+0.080$ \\
 50 & 0.567 & $0.567$ & $0.000$ \\
100 & 0.533 & $0.517$ & $-0.016$ \\
200 & 0.481 & $0.481$ & $0.000$ \\
\bottomrule
\end{tabular}
\end{center}

The fit is poor at small $N$ (residuals $> 0.1$) and good at large
$N$. This is expected: the $\eta_\infty + c/\sqrt{N}$ model captures
bulk-packing behavior but not the boundary effects that dominate at
small $N$ where only a few pieces interact strongly with the container
walls.

\gabitov{%
The $1/\sqrt{N}$ scaling is a phenomenological fit, not a theoretical
prediction. It should be reported as such. The physical interpretation
is that at large $N$, the fraction of polygons touching the container
wall (and therefore experiencing boundary effects) is proportional to
the perimeter-to-area ratio of the container, which scales as
$1/L \sim 1/\sqrt{N}$. This argument is heuristic. A rigorous bound
on $\eta^*$ as a function of $N$ for this specific polygon is an open
problem.%
}

\expcheck{%
The model predicts $\eta_\infty \approx 0.395$ as the large-$N$ limit.
This seems low: high-quality packers for convex shapes reach $\eta
\approx 0.9$, and non-convex shapes with good interlocking can exceed
that. The low observed density at $N = 200$ (0.481) compared to $N = 50$
(0.567) is more likely a consequence of under-exploration (the solver
finds feasibility but not a particularly tight packing) than a true
physical limit. The extrapolated $\eta_\infty = 0.395$ should be
treated with skepticism; it will improve as the solver quality increases
with longer wall times and more workers.%
}

% ============================================================
\section{Comparison: Predictions vs.\ Experimental Results}
\label{sec:comparison}

This section walks through each quantitative prediction made in the
companion report and checks it against observed data.

\subsection{Bracket-Plus-Bisect Runtime}

\textbf{Prediction:} $t_{\mathrm{bb}} \approx 23$ minutes per worker.

\textbf{Observed:} Median first-feasible times (Table~2 in the report)
were 5--53 seconds. These are much shorter than 23 minutes. The
reconciliation: the first-feasible time is not the end of bisection;
it is when the SA first finds a feasible configuration within any
bracket or bisection step. A single Phase B run of 35{,}000 iterations
takes $35{,}000 \times 25\;\mu\mathrm{s} = 875$ ms. At $N = 25$, a
median first-feasible time of 8 seconds corresponds to finding
feasibility around Phase B of trial~9 (each trial is about 2.6 s, and
first-feasible at 8 s $\approx$ trial 3). The 23-minute estimate is
the worst-case time to complete all bracket-plus-bisect steps, not the
time to first feasibility.

\subsection{Feasibility Rate vs.\ $N$}

\textbf{Prediction:} The configuration space expands exponentially
with $N$, so the fraction of workers reaching feasibility should
decrease with $N$.

\textbf{Observed:} 15/15 workers feasible at $N = 5$ and $N = 10$;
6/6 at $N = 25$; 4 at $N = 50$; 3 at $N = 100$; 1 at $N = 200$.

The drop is steeper than exponential scaling alone would predict.
Part of the explanation lies in the fixed iteration budget: at $N =
200$, each polygon receives $M/N = 525$ move proposals per trial,
versus 4{,}200 per polygon at $N = 25$. The SA is exploring a 600-dimensional
space ($N = 200$, 3 DOF per polygon) with the same total iteration
budget as a 75-dimensional one. A natural fix is to scale $M \propto N$,
which would multiply per-trial time linearly with $N$ but keep the
per-polygon exploration budget constant.

\subsection{Time to Final Best Configuration}

\textbf{Prediction:} The majority of wall time goes to the polish phase,
which should dominate the time-to-best statistic.

\textbf{Observed:} Median times to best at $N = 50$ and $N = 100$ were
3{,}657 s and 3{,}254 s respectively, out of 14{,}100 s total wall time.
This places the best improvement at roughly 23--26\% into the job,
with the polish phase responsible. At $N = 25$, median time to best was
205 s (1.5\% of wall time), suggesting the bisection phase found a good
solution quickly and polish provided only marginal gains. This is
consistent with $N = 25$ being a small enough problem that the SA
explores effectively within the bracket-phase budget.

% ============================================================
\section{Audit Checklist}
\label{sec:checklist}

The following checklist can be used to verify independently that every
number in this guide and the companion report is correct.

\begin{enumerate}[leftmargin=*,label=\textbf{\arabic*.}]

\item \textbf{Polygon area.} Run:
\begin{lstlisting}[style=bashstyle]
python3 -c "
V = [(0,0.8),(0.125,0.5),(0.0625,0.5),(0.2,0.25),(0.1,0.25),
     (0.35,0),(0.075,0),(0.075,-0.2),(-0.075,-0.2),(-0.075,0),
     (-0.35,0),(-0.1,0.25),(-0.2,0.25),(-0.0625,0.5),(-0.125,0.5)]
n=len(V)
A=0.5*abs(sum(V[i][0]*V[(i+1)%n][1]-V[(i+1)%n][0]*V[i][1]
              for i in range(n)))
print(f'A_p = {A}')
"
\end{lstlisting}
Expected output: \texttt{A\_p = 0.245625}.

\item \textbf{Bounding radius.} Run:
\begin{lstlisting}[style=bashstyle]
python3 -c "
import math
V = [(0,0.8),(0.125,0.5),(0.0625,0.5),(0.2,0.25),(0.1,0.25),
     (0.35,0),(0.075,0),(0.075,-0.2),(-0.075,-0.2),(-0.075,0),
     (-0.35,0),(-0.1,0.25),(-0.2,0.25),(-0.0625,0.5),(-0.125,0.5)]
print(f'r_br = {max(math.hypot(x,y) for x,y in V):.6f}')
"
\end{lstlisting}
Expected output: \texttt{r\_br = 0.800000}.

\item \textbf{Area lower bounds.} For each $N$, run:
\begin{lstlisting}[style=bashstyle]
python3 -c "
import math
Ap = 0.245625
for N in [5,10,25,50,100,200]:
    print(f'N={N:4d}  L_area={math.sqrt(N*Ap):.6f}')
"
\end{lstlisting}

\item \textbf{Grid neighbor count.} Verify equation~\eqref{eq:keta}:
\begin{lstlisting}[style=bashstyle]
python3 -c "
Ap=0.245625; rbr=0.8; h=2*rbr
k = 25*h**2/Ap
print(f'k(eta) = {k:.2f} * eta')
print(f'k(0.85) = {k*0.85:.1f}')
"
\end{lstlisting}
Expected: \texttt{260.6 * eta}, \texttt{221.5} at $\eta = 0.85$.

\item \textbf{AABB pass fraction.} Verify equation~\eqref{eq:faabb}:
\begin{lstlisting}[style=bashstyle]
python3 -c "
rbr=0.8; h=2*rbr
exclusion = (2*rbr)**2
window    = 25*h**2
print(f'AABB pass fraction = {exclusion/window:.3f}')
"
\end{lstlisting}
Expected: \texttt{0.040}.

\item \textbf{Per-move flop count.} Verify equation~\eqref{eq:cmove}:
\begin{lstlisting}[style=bashstyle]
python3 -c "
k_nbr = 222; C_aabb = 4
k_sat = 9;   C_sat  = 2400
C_move = k_nbr*C_aabb + k_sat*C_sat
print(f'C_move = {C_move}')
"
\end{lstlisting}
Expected: \texttt{22488} $\approx 22{,}500$.

\item \textbf{Time per move.}
$t_{\mathrm{move}} = 22{,}500 / 10^9 \;\mathrm{s} = 22.5\;\mu\mathrm{s}$.
Verify empirically by wrapping the inner SA loop in \texttt{clock()}
calls over 100{,}000 iterations and dividing by $10^5$.

\item \textbf{Bracket-plus-bisect time.} Verify equation~\eqref{eq:tbb}:
\begin{lstlisting}[style=bashstyle]
python3 -c "
M=105000; K_br=10; K_bi=7
steps_br=35; steps_bi=26
I_bb = (steps_br*K_br + steps_bi*K_bi)*M
t_move_us = 25
t_bb_min = I_bb*t_move_us/1e6/60
print(f'I_bb = {I_bb:.2e}')
print(f't_bb = {t_bb_min:.1f} min')
"
\end{lstlisting}
Expected: \texttt{I\_bb = 5.59e+07}, \texttt{t\_bb = 23.3 min}.

\item \textbf{Bisection precision.}
$\Delta L_{26} = \Delta L_0 / 2^{26}$. For $\Delta L_0 = 0.25$:
$0.25 / 67{,}108{,}864 = 3.72 \times 10^{-9}$. Verify with:
\texttt{print(0.25 / 2**26)}.

\item \textbf{Expected overlapping pairs (upper bound).} Verify
equation~\eqref{eq:Eov} at $N = 100$, $\eta = 0.85$:
\begin{lstlisting}[style=bashstyle]
python3 -c "
import math
N=100; eta=0.85; rbr=0.8; Ap=0.245625
E_ov = 2*(N-1)*math.pi*rbr**2*eta/Ap
print(f'E[n_ov] <= {E_ov:.1f}')
"
\end{lstlisting}
Expected: approximately $1{,}387$.

\item \textbf{Tolerance depth.} Verify equation~\eqref{eq:dtol}
for Phase A start ($T = 0.30$, $\lambda = 100$, $p = 1$):
$d_{\mathrm{tol}} = 0.30/100 = 0.003$.
For Phase B start ($T = 0.06$, $\lambda = 1000$, $p = 2$):
$d_{\mathrm{tol}} = \sqrt{0.06/1000} = 7.75 \times 10^{-3}$.

\item \textbf{Cooling rate.} Verify $\beta_A$ from
equation~\eqref{eq:cool}:
\begin{lstlisting}[style=bashstyle]
python3 -c "
T0=0.30; Tend=3e-4; M=70000
beta = (Tend/T0)**(1/M)
print(f'beta_A = {beta:.8f}')
print(f'T at 35000 iters = {T0*beta**35000:.4e}')
"
\end{lstlisting}
Expected: $\beta_A \approx 0.99990152$; midpoint temperature
$\approx 0.0095$.

\item \textbf{Worker count.} Check the SLURM job log for the line
\texttt{WORKERS: N} and confirm $N = (\text{detected CPUs}) - 2$.
On the 32-CPU nodes, this should be 30.

\item \textbf{Area efficiency in the output.} The SLURM scripts compute
area efficiency automatically on completion:
\begin{lstlisting}[style=bashstyle]
python3 -c "N=100; area=0.245625; L=6.785; \
            print(N*area/L**2)"
\end{lstlisting}
Expected for $N=100$, $L^*=6.785$: $\eta = 0.534$.

\end{enumerate}

% ============================================================
\section{Summary of All Key Constants}
\label{sec:summary}

\begin{center}
\begin{tabular}{lll}
\toprule
Quantity & Symbol & Value \\
\midrule
Polygon vertex count        & $N_V$               & 15 \\
Triangle count              & $T$                 & 13 \\
Triangle groups             & ---                 & (5, 5, 3) \\
Polygon area                & $A_p$               & 0.245625 \\
Bounding radius             & $r_{\mathrm{br}}$   & 0.80 \\
Grid cell size              & $h$                 & 1.60 \\
Search radius (cells)       & $R$                 & 2 \\
Grid window (cells)         & ---                 & $5\times 5 = 25$ \\
\midrule
Phase A iterations          & $M_A$               & 70{,}000 \\
Phase B iterations          & $M_B$               & 35{,}000 \\
Total per trial             & $M$                 & 105{,}000 \\
Phase A: $T_{\mathrm{start}}$ & ---              & 0.30 \\
Phase A: $T_{\mathrm{end}}$   & ---              & $3\times10^{-4}$ \\
Phase B: $T_{\mathrm{start}}$ & ---              & 0.06 \\
Phase B: $T_{\mathrm{end}}$   & ---              & $10^{-6}$ \\
Adapt window               & ---                 & 1{,}500 iterations \\
\midrule
Neighbour count (coeff.)   & $k(\eta)/\eta$      & $\approx 261$ \\
AABB pass fraction         & $f_{\mathrm{AABB}}$ & 4\% \\
SAT cost per pair          & $C_{\mathrm{SAT}}$  & $\approx 2{,}400$ flops \\
Per-move cost ($\eta=0.85$)& $C_{\mathrm{move}}$ & $\approx 22{,}500$ flops \\
Time per move              & $t_{\mathrm{move}}$ & $\approx 25\;\mu\mathrm{s}$ \\
\midrule
Bracket-bisect iterations  & $I_{\mathrm{bb}}$   & $5.6\times10^7$ \\
Bracket-bisect time (1 wkr)& $t_{\mathrm{bb}}$   & $\approx 23$ min \\
Bisection steps            & $k$                 & 26 \\
Bisection precision        & $\Delta L_{26}/\Delta L_0$  & $1.5\times10^{-8}$ \\
Feasibility tolerance      & ---                 & $10^{-10}$ \\
\midrule
Nodes (5-node jobs)        & ---                 & 5 \\
CPUs per node              & ---                 & 32 \\
Workers per node           & $W$                 & 30 \\
Total concurrent workers   & $W_{\mathrm{total}}$& 150 \\
Runs queued per node       & ---                 & 64 \\
Wall time ($N \leq 100$)   & ---                 & 4 h = 14{,}400 s \\
Wall time ($N = 200$)      & ---                 & 8 h = 28{,}800 s \\
Effective time limit       & ---                 & Wall $-$ 300 s \\
\bottomrule
\end{tabular}
\end{center}

% ============================================================
\bibliographystyle{plain}
\begin{thebibliography}{5}

\bibitem{kirkpatrick1983}
S.~Kirkpatrick, C.~D. Gelatt, and M.~P. Vecchi.
Optimization by simulated annealing.
\textit{Science}, 220(4598):671--680, 1983.

\bibitem{leung2003}
T.~W. Leung, C.~K. Chan, and M.~D. Troutt.
Application of a mixed simulated annealing-genetic algorithm heuristic
for the two-dimensional orthogonal packing problem.
\textit{European Journal of Operational Research}, 145(3):530--542, 2003.

\bibitem{Hajek1988}
B.~Haj\'ek.
Cooling schedules for optimal annealing.
\textit{Mathematics of Operations Research}, 13(2):311--329, 1988.

\bibitem{bennell2008}
J.~A. Bennell and J.~F. Oliveira.
The geometry of nesting problems: A tutorial.
\textit{European Journal of Operational Research}, 184(2):397--415, 2008.

\end{thebibliography}

\end{document}